\scanpage{24}
% L1–4
Inversement, soit $D_{*}$ un ens. ordonné satisfaisant \textcircled{a} ci-dessus. Donnons nous sur $D_0$ (ens des él. minimaux de $D_*$) une structure d'ordre (pas seulement de préordre) \inserted{notée $x\unsure{\preceq|\leq|\preccurlyeq}y$}, \deleted{\unsure{traduisant|induisant} \unsure{ce que}}
% L5–8
\inserted{satisfaisant \textcircled{b} ci-dessus et de plus}
\deleted{\unsure{signifiant que|vérifiant que}} \inserted{\textcircled{c} Elle \deleted{\unsure{sur}} induit sur} les $I_d$ \deleted{$(d\in D_*)$} $=\{x\in D_0\mid x\leq d\}$ un ordre total. \deleted{qui} \inserted{Ce dernier} est évidemment fonctoriel \deleted{\unsure{ses ordonnés|pour ordonné} (pour l'ordre induit par}
% L9–12
en $d$, de sorte que $D_*$ définit un foncteur \unsure{croissant|covariant} $\underline{D_*}\longrightarrow\unsure{\mathbb{\Delta}|\Delta}$, qui est \uline{fibrant}, à fibres discrètes -- donc associé à un foncteur $D_*:\mathbb{\Delta}^{\circ}\longrightarrow(\mathrm{Ens})$, et $D_0=D_*(\Delta_{\unsure{n|0}})$
% L13–15
(plus gén. $D_n=D_*(\Delta_n)$ est l'ens des $d\in D_*$ tels que $\mathrm{Card}(I_d)=n+1$ \unsure{$\}$}, i.e. $I_d\simeq\Delta_n$), et $D_*$ satisfait évidemment \textcircled{1}.

% L16–23 left
\diagram{\begin{tikzcd}[ampersand replacement=\&,row sep=large]
\begin{gathered}\text{Ens. quasi-simpliciaux }\text{\inserted{$D_*$}}\text{ satisfaisant}\\
\text{\textcircled{1} et \textcircled{2}}\end{gathered}\arrow[d,shift right=5]\\
\begin{gathered}\text{Ens ordonnés }\text{\inserted{$D_*$}}\text{ satisfaisant \textcircled{a} et}\\
\text{munis d'une relation d'ordre}\\
\text{sur }D_0=\text{ ens des él minimaux}\\
\text{de }D_*\text{, satisfaisant les conditions}\\
\text{\textcircled{b} et \textcircled{c}}\end{gathered}\arrow[u,shift right=5]
\end{tikzcd}}
% L16–20 right
Soient alors $x_0,y_0\in D_0$ tels que $x_0\prec_{*}y_0$ i.e. $\exists z_1\in\unsure{\mathbb D|D}_1$, $x_0=\delta_0(z_1),y_0=\delta_1(z_1)$ avec $z_1\in D_1$. Donc $I_{z_1}=\{x_0,y_0\}$, et par définition de $\delta_0,\delta_1$, $x_0$ est
% L21–24 right
le $\mathrm{1}^{\mathrm{er}}$ él de $I_{z_1}$ pour l'ordre \inserted{(total)} induit par $D_0$ $y_0$ le deuxième, i.e. on a $x_0\prec y_0$. Inversement,
% L25–27
supposons $x_0\prec y_0$, \deleted{donc} $\exists\unsure{z|d}\in D_*$ tel que $x_0<z,y_0<z$ -- on aura $z\in D_n$ avec $n\geq1$, $x_0,y_0$ él. distincts de $I_z$, \deleted{donc} on a $x_0\prec y_0$ pour
% archive sheet number: 23
